\documentclass[12pt,a4paper]{article}
\usepackage[utf8]{inputenc}
\usepackage[T1]{fontenc}
\usepackage[polish]{babel}
\usepackage{geometry}
\usepackage{enumitem}
\usepackage{titlesec}
\usepackage{parskip}
\usepackage{fancyhdr}
\usepackage{xcolor}
\usepackage{mdframed}
\usepackage{microtype}
\usepackage{booktabs}
\usepackage{hyperref}
\usepackage{multirow}
\usepackage{array}
\usepackage{longtable}
\usepackage{tabularx}

\geometry{margin=2.5cm}
\hypersetup{colorlinks=true, linkcolor=black, urlcolor=blue}

\pagestyle{fancy}
\fancyhf{}
\rhead{SDLC-LLM Benchmark}
\lhead{Zadanie R2 -- Refinement}
\rfoot{\thepage}

\titleformat{\section}{\large\bfseries}{}{0em}{}[\titlerule]
\titleformat{\subsection}{\normalsize\bfseries}{}{0em}{}
\titleformat{\subsubsection}{\normalsize\itshape}{}{0em}{}

\definecolor{metabg}{gray}{0.93}
\definecolor{promptbg}{RGB}{235,245,255}
\definecolor{claudecolor}{RGB}{210,105,30}
\definecolor{score3}{RGB}{220,240,220}
\definecolor{score2}{RGB}{255,243,205}
\definecolor{score1}{RGB}{255,220,200}
\definecolor{score0}{RGB}{255,200,200}
\definecolor{tpcolor}{RGB}{220,240,220}
\definecolor{fpcolor}{RGB}{255,220,200}
\definecolor{fncolor}{RGB}{255,243,205}

\newcommand{\scorebox}[1]{%
  \ifnum#1=3\colorbox{score3}{\textbf{3}}%
  \else\ifnum#1=2\colorbox{score2}{\textbf{2}}%
  \else\ifnum#1=1\colorbox{score1}{\textbf{1}}%
  \else\colorbox{score0}{\textbf{0}}\fi\fi\fi
}
\newcommand{\tp}[1]{\colorbox{tpcolor}{\small #1}}
\newcommand{\fp}[1]{\colorbox{fpcolor}{\small #1}}
\newcommand{\fn}[1]{\colorbox{fncolor}{\small #1}}

\begin{document}

% ─── STRONA TYTUŁOWA ────────────────────────────────────────────────
\begin{center}
  {\LARGE\bfseries Requirements Artifact}\\[1.5em]
  {\large \textbf{Case study:} Appointment Booking System}\\[0.3em]
  {\large \textbf{Model:} \texttt{gpt-5.4}}\\[0.3em]
  {\large \textbf{Tryb:} LLM-only}\\[0.3em]
  {\large \textbf{Wejście:} Gold ver2 (challenging -- celowe niejednoznaczności)}
\end{center}

\vspace{1em}\hrule\vspace{1.5em}

% ─── PROMPT ─────────────────────────────────────────────────────────
\section{Prompt}

\begin{mdframed}[backgroundcolor=metabg, linewidth=0pt, innerleftmargin=10pt, innerrightmargin=10pt, innertopmargin=8pt, innerbottommargin=8pt]
\textbf{System prompt:}\\
\small Jesteś analitykiem wymagań. Odpowiadaj WYŁĄCZNIE poprawnym JSON. Żadnego Markdown ani komentarzy.
\end{mdframed}

\vspace{0.5em}
\begin{mdframed}[backgroundcolor=promptbg, linewidth=0pt, innerleftmargin=10pt, innerrightmargin=10pt, innertopmargin=8pt, innerbottommargin=8pt]
\textbf{User prompt (3-krokowy):}\\[0.3em]
\small
\texttt{[KROK 1]} Na podstawie opisu systemu wygeneruj \texttt{system\_goal}, role i wymagania funkcjonalne.\\
\texttt{[KROK 2]} Na podstawie opisu i wygenerowanych FR wygeneruj NFR i user stories.\\
\texttt{[KROK 3]} Na podstawie user stories wygeneruj acceptance criteria w formacie Given/When/Then.\\[0.3em]
\textbf{Opis systemu (wejście):}\\
System umożliwia użytkownikom rezerwację wizyt u specjalistów. Użytkownicy mogą przeglądać dostępne terminy, dokonywać rezerwacji oraz je anulować. Specjaliści zarządzają swoim grafikiem. System powinien zapobiegać konfliktom rezerwacji, ale w niektórych przypadkach dopuszczalne są wyjątki. Użytkownik może anulować wizytę w uzasadnionych przypadkach, jednak anulowanie może być ograniczone czasowo. Specjalista może modyfikować swój grafik, także dla już istniejących rezerwacji. System powinien obsługiwać różne role użytkowników i być skalowalny.
\end{mdframed}

% ─── CEL SYSTEMU ────────────────────────────────────────────────────
\section{Cel systemu}

System umożliwia użytkownikom rezerwację, przeglądanie i anulowanie wizyt u specjalistów oraz wspiera specjalistów w zarządzaniu grafikiem. Celem systemu jest bezpieczna i skalowalna obsługa terminów, rezerwacji i zmian harmonogramu przy uwzględnieniu ról użytkowników, ograniczeń czasowych oraz wyjątków od konfliktów rezerwacji.

% ─── ROLE ───────────────────────────────────────────────────────────
\section{Role}

\begin{itemize}[leftmargin=1.5em]
  \item \textbf{Użytkownik} -- przeglądanie dostępnych terminów, dokonywanie rezerwacji wizyt, anulowanie własnych rezerwacji w dozwolonych warunkach
  \item \textbf{Specjalista} -- zarządzanie własnym grafikiem, modyfikowanie dostępnych terminów, modyfikowanie grafiku także dla istniejących rezerwacji, przeglądanie rezerwacji dotyczących własnego grafiku
  \item \textbf{Administrator} -- zarządzanie rolami użytkowników, konfiguracja zasad rezerwacji i anulowania, obsługa wyjątków od konfliktów rezerwacji, nadzór nad skalowalnością i poprawnym działaniem systemu
\end{itemize}

% ─── WYMAGANIA FUNKCJONALNE ─────────────────────────────────────────
\section{Wymagania funkcjonalne}

\begin{description}[leftmargin=1.8cm, labelwidth=1.3cm, labelsep=0.5cm, style=nextline]
  \item[FR1] System musi umożliwiać użytkownikowi przeglądanie dostępnych terminów wizyt u specjalistów.
  \item[FR2] System musi umożliwiać użytkownikowi dokonanie rezerwacji wybranego dostępnego terminu.
  \item[FR3] System musi umożliwiać użytkownikowi anulowanie własnej rezerwacji, jeśli spełnione są obowiązujące warunki anulowania.
  \item[FR4] System musi umożliwiać definiowanie i egzekwowanie ograniczeń czasowych dotyczących anulowania wizyt.
  \item[FR5] System musi zapobiegać konfliktom rezerwacji dla tego samego specjalisty i terminu.
  \item[FR6] System musi umożliwiać obsługę wyjątków od standardowych reguł konfliktu rezerwacji w określonych przypadkach.
  \item[FR7] System musi umożliwiać specjaliście zarządzanie własnym grafikiem, w tym dodawanie, usuwanie i zmianę dostępnych terminów.
  \item[FR8] System musi umożliwiać specjaliście modyfikowanie grafiku także wtedy, gdy istnieją już powiązane rezerwacje.
  \item[FR9] System musi obsługiwać różne role użytkowników oraz przypisane im uprawnienia.
  \item[FR10] System musi być skalowalny i wspierać rosnącą liczbę użytkowników, specjalistów, grafików i rezerwacji.
\end{description}

% ─── WYMAGANIA NIEFUNKCJONALNE ──────────────────────────────────────
\section{Wymagania niefunkcjonalne}

\begin{description}[leftmargin=1.8cm, labelwidth=1.3cm, labelsep=0.5cm, style=nextline]
  \item[NFR1] System powinien zapewniać wysoką dostępność funkcji przeglądania terminów, rezerwacji, anulowania i zarządzania grafikiem.
\item[NFR2] System powinien zapewniać spójność danych i odporność na współbieżne próby rezerwacji tego samego terminu.
\item[NFR3] System powinien skalować się wraz ze wzrostem liczby użytkowników, specjalistów, grafików i rezerwacji bez istotnego pogorszenia wydajności.
\item[NFR4] System powinien egzekwować kontrolę dostępu opartą na rolach użytkowników zgodnie z przypisanymi uprawnieniami.
\item[NFR5] System powinien zapewniać bezpieczeństwo danych użytkowników i rezerwacji, w tym ochronę przed nieautoryzowanym dostępem.
\item[NFR6] System powinien rejestrować operacje związane z rezerwacjami, anulowaniami, wyjątkami konfliktów oraz zmianami grafiku w celu audytu i rozliczalności.
\item[NFR7] System powinien być projektowany w sposób umożliwiający łatwą konfigurację zasad anulowania, wyjątków od konfliktów oraz uprawnień ról.
\item[NFR8] System powinien zapewniać czytelny i intuicyjny interfejs dla użytkowników, specjalistów i administratorów.
\end{description}

% ─── USER STORIES ───────────────────────────────────────────────────
\section{User stories}

\begin{description}[leftmargin=1.8cm, labelwidth=1.3cm, labelsep=0.5cm, style=nextline]
  \item[US1] Jako użytkownik chcę przeglądać dostępne terminy u specjalistów, aby wybrać dogodny termin wizyty.
\item[US2] Jako użytkownik chcę zarezerwować dostępny termin wizyty, aby umówić spotkanie ze specjalistą.
\item[US3] Jako użytkownik chcę anulować własną rezerwację w dozwolonych warunkach, aby móc zrezygnować z wizyty, gdy jest to uzasadnione.
\item[US4] Jako użytkownik chcę otrzymać informację o ograniczeniach anulowania, aby wiedzieć, czy mogę jeszcze odwołać wizytę.
\item[US5] Jako specjalista chcę zarządzać własnym grafikiem, aby udostępniać i aktualizować terminy wizyt.
\item[US6] Jako specjalista chcę modyfikować grafik także dla istniejących rezerwacji, aby dostosować harmonogram do zmian organizacyjnych.
\item[US7] Jako specjalista chcę przeglądać rezerwacje dotyczące mojego grafiku, aby zarządzać zaplanowanymi wizytami.
\item[US8] Jako administrator chcę zarządzać rolami i uprawnieniami użytkowników, aby kontrolować dostęp do funkcji systemu.
\item[US9] Jako administrator chcę konfigurować zasady rezerwacji i anulowania, aby dostosować działanie systemu do polityki organizacji.
\item[US10] Jako administrator chcę obsługiwać wyjątki od konfliktów rezerwacji, aby umożliwiać niestandardowe przypadki zgodnie z obowiązującymi zasadami.
\end{description}

% ─── ACCEPTANCE CRITERIA ────────────────────────────────────────────
\section{Acceptance criteria}

\textbf{AC1} -- Przeglądanie dostępnych terminów
\begin{itemize}[itemsep=2pt]
\item \textbf{Given} użytkownik jest na ekranie przeglądania terminów i w systemie istnieją opublikowane wolne terminy specjalistów
\item \textbf{When} użytkownik wyświetla listę dostępnych terminów
\item \textbf{Then} system pokazuje dostępne terminy wraz z informacją o specjaliście, dacie i godzinie wizyty
\end{itemize}

\textbf{AC2} -- Filtrowanie terminów
\begin{itemize}[itemsep=2pt]
\item \textbf{Given} użytkownik przegląda dostępne terminy specjalistów
\item \textbf{When} użytkownik wybiera kryteria wyszukiwania lub filtrowania terminów
\item \textbf{Then} system ogranicza widok do terminów spełniających wybrane kryteria
\end{itemize}

\textbf{AC3} -- Potwierdzenie rezerwacji
\begin{itemize}[itemsep=2pt]
\item \textbf{Given} użytkownik jest zalogowany i wybrany termin jest nadal dostępny do rezerwacji
\item \textbf{When} użytkownik potwierdza rezerwację terminu wizyty
\item \textbf{Then} system tworzy rezerwację i oznacza termin jako zarezerwowany
\end{itemize}

\textbf{AC4} -- Odrzucenie rezerwacji zajętego terminu
\begin{itemize}[itemsep=2pt]
\item \textbf{Given} użytkownik próbuje zarezerwować termin, który został już zajęty lub przestał być dostępny
\item \textbf{When} użytkownik wysyła żądanie rezerwacji
\item \textbf{Then} system odrzuca rezerwację i informuje użytkownika o braku dostępności terminu
\end{itemize}

\textbf{AC5} -- Anulowanie rezerwacji zgodnie z warunkami
\begin{itemize}[itemsep=2pt]
\item \textbf{Given} użytkownik ma własną aktywną rezerwację i spełnione są skonfigurowane warunki anulowania
\item \textbf{When} użytkownik anuluje rezerwację
\item \textbf{Then} system usuwa rezerwację lub oznacza ją jako anulowaną oraz ponownie udostępnia termin zgodnie z konfiguracją
\end{itemize}

\textbf{AC6} -- Blokada anulowania rezerwacji
\begin{itemize}[itemsep=2pt]
\item \textbf{Given} użytkownik ma własną rezerwację, ale nie są spełnione skonfigurowane warunki anulowania
\item \textbf{When} użytkownik próbuje anulować rezerwację
\item \textbf{Then} system blokuje anulowanie i wyświetla informację o obowiązujących ograniczeniach
\end{itemize}

\textbf{AC7} -- Wyświetlanie warunków anulowania
\begin{itemize}[itemsep=2pt]
\item \textbf{Given} użytkownik przegląda szczegóły własnej rezerwacji lub rozpoczyna proces anulowania
\item \textbf{When} system ocenia zasady anulowania dla tej rezerwacji
\item \textbf{Then} system pokazuje, czy anulowanie jest możliwe oraz jakie ograniczenia lub terminy obowiązują
\end{itemize}

\textbf{AC8} -- Dodawanie terminu przez specjalistę
\begin{itemize}[itemsep=2pt]
\item \textbf{Given} specjalista jest zalogowany i posiada uprawnienia do zarządzania własnym grafikiem
\item \textbf{When} specjalista dodaje nowy termin dostępności do swojego grafiku
\item \textbf{Then} system zapisuje termin i udostępnia go użytkownikom jako możliwy do rezerwacji, jeśli nie narusza reguł konfliktów
\end{itemize}

\textbf{AC9} -- Modyfikacja wolnego terminu przez specjalistę
\begin{itemize}[itemsep=2pt]
\item \textbf{Given} specjalista jest zalogowany i posiada istniejące terminy w swoim grafiku bez aktywnych rezerwacji
\item \textbf{When} specjalista modyfikuje lub usuwa wybrany termin
\item \textbf{Then} system aktualizuje grafik zgodnie ze zmianą, jeśli operacja nie narusza skonfigurowanych zasad
\end{itemize}

\textbf{AC10} -- Modyfikacja zarezerwowanego terminu przez specjalistę
\begin{itemize}[itemsep=2pt]
\item \textbf{Given} specjalista jest zalogowany i wybrany termin w jego grafiku ma istniejącą rezerwację
\item \textbf{When} specjalista modyfikuje ten termin w zakresie dozwolonym przez system
\item \textbf{Then} system zapisuje zmianę zgodnie z obowiązującymi zasadami oraz odpowiednio aktualizuje powiązaną rezerwację lub jej status
\end{itemize}

\textbf{AC11} -- Przegląd rezerwacji przez specjalistę
\begin{itemize}[itemsep=2pt]
\item \textbf{Given} specjalista jest zalogowany i posiada terminy z przypisanymi rezerwacjami
\item \textbf{When} specjalista otwiera widok rezerwacji dotyczących swojego grafiku
\item \textbf{Then} system pokazuje listę rezerwacji powiązanych wyłącznie z grafikiem tego specjalisty
\end{itemize}

\textbf{AC12} -- Zarządzanie uprawnieniami przez administratora
\begin{itemize}[itemsep=2pt]
\item \textbf{Given} administrator jest zalogowany i posiada uprawnienia do zarządzania rolami oraz uprawnieniami
\item \textbf{When} administrator przypisuje, zmienia lub odbiera rolę albo uprawnienie użytkownikowi
\item \textbf{Then} system zapisuje zmianę i stosuje nowe uprawnienia przy kolejnych operacjach użytkownika
\end{itemize}

\textbf{AC13} -- Konfiguracja zasad systemowych przez administratora
\begin{itemize}[itemsep=2pt]
\item \textbf{Given} administrator jest zalogowany i posiada uprawnienia do konfiguracji zasad systemowych
\item \textbf{When} administrator definiuje lub zmienia zasady rezerwacji oraz anulowania
\item \textbf{Then} system zapisuje konfigurację i stosuje ją do nowych prób rezerwacji oraz anulowania
\end{itemize}

\textbf{AC14} -- Obsługa wyjątków przez administratora
\begin{itemize}[itemsep=2pt]
\item \textbf{Given} występuje konflikt rezerwacji lub konflikt wynikający z reguł harmonogramu oraz administrator ma odpowiednie uprawnienia
\item \textbf{When} administrator zatwierdza wyjątek zgodnie z obowiązującymi zasadami organizacji
\item \textbf{Then} system pozwala na realizację operacji mimo standardowego konfliktu i zapisuje informację o zastosowanym wyjątku
\end{itemize}


% ─── PYTANIA UZNANE ZA ROZSTRZYGNIĘTE ──────────────────────────────
\section{Pytania uznane za rozstrzygnięte}

W tej wersji artefaktu model przyjął następujące założenia:

\begin{itemize}[leftmargin=1.5em]
  \item polityka anulowania wizyt jest możliwa do zdefiniowana przez administratora systemu
  \item brak sztywnej długości slotu czasowego wizyt
  \item jedna rezerwacja dotyczy jednego specjalisty i jednego slotu
   \item mechanizmy współbieżności nie są określone
   \item brak stref czasowych
    \item brak refundów
\end{itemize}

\newpage

% ─── METRYKI VS GOLD ────────────────────────────────────────────────
\section{Metryki porównawcze względem gold artifact}

\begin{mdframed}[backgroundcolor=metabg, linewidth=0pt, innerleftmargin=10pt, innerrightmargin=10pt, innertopmargin=8pt, innerbottommargin=8pt]
\small
\textbf{Legenda:} \tp{TP} = pokrywa element gold \quad \fp{FP} = nadmiarowe, nie ma w gold \quad \fn{FN} = brakujące względem gold\\[0.3em]
Metryki liczone na poziomie wymagań semantycznych (nie IDs), tzn.\ porównywana jest \textit{treść} wymagania z gold, nie jego numer.
\end{mdframed}

\vspace{0.8em}

\subsection{Wymagania funkcjonalne}

\begin{tabular}{p{2.5cm} p{7cm} l}
\toprule
\textbf{ID Gpt} & \textbf{Mapowanie na gold} & \textbf{Status} \\
\midrule
FR1 & Gold FR1 -- przeglądanie terminów & \tp{TP} \\
FR2 & Gold FR2 -- rezerwacja dostępnego terminu & \tp{TP} \\
FR3 & Gold FR5 -- anulowanie rezerwacji ze spełnionym warunkami & \tp{TP} \\
FR4 & Gold FR5 -- wprowadzenie ograniczeń czasowych odnoście anulowania wizyt & \tp{TP} \\
FR5 & Gold FR4 -- brak podwójnej rezerwacji & \tp{TP} \\
FR6 & Gold FR4 -- obsługa wyjątków odnoście konfliktu rezerwacji & \tp{TP} \\
FR7 & Gold FR7 -- zarządzanie grafikiem specjalisty & \tp{TP} \\
FR8 & Gold FR7 -- modyfikacja terminu w przypadku aktywnych rezerwacji & \tp{TP} \\
FR9 & Gold FR9 -- dostęp do danych & \tp{TP} \\
FR10 & brak w Gold -- skalowalność systemu & \fp{FP} \\

\midrule
\multicolumn{3}{l}{\textit{Brakujące z gold:}} \\
-- & Gold FR3 -- limit 3 aktywnych rezerwacji per user & \fn{FN} \\
-- & Gold FR6 -- zwolnienie terminu / BLOCKED & \fn{FN} \\
-- & Gold FR8 -- statusy BLOCKED, niejasne przejścia stanów & \fn{FN} \\
-- & Gold FR10 -- zakaz nakładających się rezerwacji & \fn{FN} \\
-- & Gold FR11 -- blokady slotów (BLOCKED) & \fn{FN} \\
-- & Gold FR12 -- historia operacji & \fn{FN} \\
-- & Gold FR13 -- równoczesne operacje & \fn{FN} \\
\bottomrule
\end{tabular}

\vspace{0.5em}
\begin{tabular}{lrrrr}
\toprule
& \textbf{TP} & \textbf{FP} & \textbf{FN} & \\
\midrule
FR & 9 & 1 & 7 & \\
\midrule
\textbf{Precision} & \multicolumn{4}{l}{$9 / (9+1) = \mathbf{0{,}90}$} \\
\textbf{Recall}    & \multicolumn{4}{l}{$9 / (9+7) = \mathbf{0{,}56}$} \\
\textbf{F1}        & \multicolumn{4}{l}{$2 \times 0{,}90 \times 0{,}56 / (0{,}90+0{,}56) = \mathbf{0{,}69}$} \\
\bottomrule
\end{tabular}

\vspace{1em}
\subsection{Wymagania niefunkcjonalne}

\begin{tabular}{p{2.5cm} p{7cm} l}
\toprule
\textbf{ID Gpt} & \textbf{Mapowanie na gold} & \textbf{Status} \\
\midrule
NFR2 & Gold NFR1 -- spójność danych & \tp{TP} \\
NFR3 (skalowanie)    & Gold NFR5 -- skalowalność & \tp{TP} \\
NFR4 (autoryzacja)   & Gold NFR3 -- bezpieczeństwo & \tp{TP} \\
NFR5 (bezpieczeństwo danych)        & Gold NFR3 -- bezpieczeństwo & \tp{TP} \\
NFR6 (log audytowy)  & Gold NFR4 -- audyt & \tp{TP} \\
NFR1, NFR7, NFR8 & brak bezpośredniego mapowania na gold & \fp{FP} \\
\midrule
-- & Gold NFR2 -- wydajność operacji & \fn{FN} \\
\bottomrule
\end{tabular}

\vspace{0.5em}
\begin{tabular}{lrrrr}
\toprule
& \textbf{TP} & \textbf{FP} & \textbf{FN} & \\
\midrule
NFR & 5 & 3 & 1 & \\
\midrule
\textbf{Precision} & \multicolumn{4}{l}{$5 / (5+3) = \mathbf{0{,}63}$} \\
\textbf{Recall}    & \multicolumn{4}{l}{$5 / (5+1) = \mathbf{0{,}83}$} \\
\textbf{F1}        & \multicolumn{4}{l}{$2 \times 0{,}63 \times 0{,}83 / (0{,}63+0{,}83) = \mathbf{0{,}72}$} \\
\bottomrule
\end{tabular}

\vspace{1em}
\subsection{Acceptance criteria}

\begin{tabular}{p{2.5cm} p{7cm} l}
\toprule
\textbf{ID Gpt} & \textbf{Mapowanie na gold} & \textbf{Status} \\
\midrule
AC3 & Gold AC1 -- rezerwacja wizyty & \tp{TP} \\
AC4 & Gold AC3 -- odrzucenie poodwójnej rezerwacji  & \tp{TP} \\
AC5--AC6 & Gold AC4 -- anulowanie wizyty & \tp{TP} \\
AC14 & Gold AC6 -- konflikt rezerwacji  & \tp{TP} \\
AC1, AC2, AC7--AC13 & brak mapowania na gold AC & \fp{FP} \\
\midrule
-- & Gold AC2 -- limit 3 rezerwacji & \fn{FN} \\
-- & Gold AC5 -- anulowanie graniczne (dokładnie 24h) & \fn{FN} \\
\bottomrule
\end{tabular}

\vspace{0.5em}
\begin{tabular}{lrrrr}
\toprule
& \textbf{TP} & \textbf{FP} & \textbf{FN} & \\
\midrule
AC & 5 & 9 & 2 & \\
\midrule
\textbf{Precision} & \multicolumn{4}{l}{$5 / (5+9) = \mathbf{0{,}36}$} \\
\textbf{Recall}    & \multicolumn{4}{l}{$5 / (5+2) = \mathbf{0{,}71}$} \\
\textbf{F1}        & \multicolumn{4}{l}{$2 \times 0{,}36 \times 0{,}71 / (0{,}36+0{,}71) = \mathbf{0{,}48}$} \\
\bottomrule
\end{tabular}

\vspace{0.5em}
\begin{mdframed}[backgroundcolor=metabg, linewidth=0pt, innerleftmargin=10pt, innerrightmargin=10pt, innertopmargin=6pt, innerbottommargin=6pt]
\small\textbf{Komentarz do AC:} Niska precision (0.36) wynika z tego, że model wygenerował łącznie 14 AC pokrywające zaawansowane scenariusze (konfiguracja admina, zarządzanie rezerwacjami, funkcjonalności specjalisty), których gold ver2 nie wymaga -- gold jest celowo minimalistyczny. Natomiast wartość recall 0.71 oznacza, że model pominął 2 kluczowe AC: limit rezerwacji oraz przypadek graniczny anulowania rezerwacji dokładnie 24h przed terminem.
\end{mdframed}
\newpage
% ─── METADANE I OCENA ───────────────────────────────────────────────
\section{Metadane i ocena}
\begin{table}[h]
% Zmieniamy tabular na tabularx, ustawiamy szerokość na \textwidth 
% i zmieniamy typ drugiej kolumny z 'l' na 'X'
\begin{tabularx}{\textwidth}{l X}
\toprule
\textbf{Pole} & \textbf{Wartość} \\
\midrule
Model & \texttt{gpt-5.4} \\
Tryb & LLM-only \\
Case study & appointment-booking (ver2 challenging) \\
Czas łącznie & 76,24s \\
\midrule
\multicolumn{2}{l}{\textit{Ocena (0--3)}} \\
\midrule
Correctness (M1)     & \scorebox{2} \quad merytorycznie poprawny, ale występuje nadmierne rozdrobnienia funkcjonalności oraz AC\\
Completeness (M2)    & \scorebox{2} \quad pokrywa główne FR,  ale brakuje m.in limitu rezerwacji, statusów rezerwacji, ograniczeń równoczesnych rezerwacj \\
Consistency (M3)     & \scorebox{3} \quad wewnętrznie spójny, brak sprzeczności \\
Clarity (M4)         & \scorebox{3} \quad czytelna struktura, AC w formacie Given/When/Then \\
Maintainability (M5) & \scorebox{2} \quad użyteczny jako wejście dla Architecture Team po uzupełnieniu brakujących FR \\
\midrule
\multicolumn{2}{l}{\textit{Metryki vs gold}} \\
\midrule
FR Precision / Recall / F1   & 0.90 / 0.56 / 0.69 \\
NFR Precision / Recall / F1  & 0.63 / 0.83 / 0.72 \\
AC Precision / Recall / F1   & 0.36 / 0.71 / 0.48 \\
\bottomrule
\end{tabularx}
\end{table}

\end{document}